\subsection{Методы оценки размерности многомерных данных}\label{sec:lit_dimensionality}
Как показано в секции~\ref{sec:population_coding}, нейронная активность сосредоточена вблизи низкоразмерного нелинейного многообразия. Количественная оценка размерности этого многообразия является важной задачей, для решения которой разработан ряд методов.

Степень, в которой нейронная активность может быть сжата в низкоразмерное пространство, зависит от синхронизации активности отдельных клеток. Для количественной характеристики этого свойства используется так называемая \emph{эффективная размерность} (ЭР) — число независимых переменных, которые создавали бы такую же структуру ковариации и, таким образом, достаточны для описания многомерных данных \cite{edim}.

В отсутствие какой-либо синхронизации эффективная размерность равна числу исходных переменных, а когда временные ряды полностью совпадают, она становится равной 1.

Следует отметить, что эффективная размерность всегда больше или равна внутренней (т.е. истинной размерности многообразия, на котором лежит многомерная активность). Разница между этими величинами заключается в том, что эффективная размерность считает это многообразие линейным (т.е. не учитывает его возможную кривизну), а внутренняя пытается описать его с помощью как можно меньшего числа нелинейных коллективных переменных.

Недавние данные указывают на то, что степень нелинейности многообразия активности может варьироваться в зависимости от конкретной изучаемой области мозга, характера выполняемой задачи и т.д. \cite{De2023, Fortunato2023}. Следовательно, количественная характеристика нелинейности структуры нейронной активности представляет большой теоретический интерес. Однако вычисление внутренней размерности, как правило, является сложной задачей, часто требующей существенных вычислительных ресурсов и специфических предположений о природе данных. В отличие от этого, эффективная размерность может быть легко вычислена на основе корреляционной матрицы многомерного временного ряда и обеспечивает верхнюю границу для внутренней размерности.

\textit{Внутренняя размерность} (ВР) — это количество коллективных мод (осей в некотором нелинейном пространстве), которое отражает природу информации, закодированной в коллективной активности. Это отличает её от различных мер «размерности вложения» — приблизительного числа измерений, исследуемых нейронным многообразием в евклидовом пространстве \cite{Jazayeri2021}.
